\sscp{Watubela}{11-03-01}
Watubela is an important witness with *q/*k > \it{k} in nearly all cases.
Other STB languages usually have *q > {\0},
but the reflexes of *k are distinct in Bati (*k = \it{k}),
Geser (*k = \it{k}),
Gorom (*k > \it{ʔ}),
and Kowiai (*k > {\0}, \it{ʔ}).
This means we must reconstruct *q = **q to Proto-Eastern Islands.
Examples of *q and *k are given in \trf{11-19} and \trf{11-20} respectively.
Kesui, which is closely related to Watubela,
has *q/*k > {\0}, \it{h}.
An example of *k > \it{h} in Kesui is *wakaR > ‹ai áha›
‘root’ and an example of *q > \it{h} is *zaqat > ‹ráhat› ‘bad’.

\begin{table}\tcp{\tsc{Eastern Islands} *q}{11-19}
\begin{adjustbox}{width=\textwidth}
	\st{2pt}
	\begin{tabular}{llllllll}\lsptoprule
local gl.	&	egg	&	head	&	thatch	&	sand	&	black	&	old	&	low tide\\
PMP	&	*\tbr{q}atəluR	&	*\tbr{q}ulu	&	*\tbr{q}atəp	&	*\tbr{q}ənay	&	*ma-\tbr{q}itəm	&	*ma-tu\tbr{q}ah	&	*ma-\tbr{q}əti\\	\midrule
Bati	&	‹tolor›	&	‹lunini›	&		&		&	‹miatan›	&		&	\\
Geser	&	{\hr}tolu	&	{\hr}ilú	&	\cg{(barem)}	&	{\hr}ena	&	{\hr}metan	&	{\hr}matuu	&	\cg{(tora)}\\
Gorom	&	{\hr}tolur	&	{\hr}ilú	&	\cg{(barem)}	&	{\hr}ena	&	{\hr}metan	&	\cg{(mataʔena)}	&	{\hr}moti\\
Watubela	&	{\hr}\tbr{k}atlu	&	{\hr}\tbr{k}ulu	&	{\hr}\tbr{k}ataf	&	{\hr}\tbr{k}ena wo	&	{\hr}ma\tbr{k}etan	&	{\hr}mtu\tbr{k}a	&	{\hr}m\tbr{k}oti\\
Kesui	&	‹atulú›	&	‹alúda›	&		&		&	‹meten›	&		&	\\
Kowiai	&	{\hr}toron	&	\cg{(umun)}	&	\cg{(barem)}	&	{\hr}ena	&	{\hr}maetan	&	\cg{(na-ɸameteʔena)}	&	{\hr}mot\\
		\lspbottomrule
	\end{tabular}
\end{adjustbox}
\end{table}

\begin{table}\tcp{\tsc{Eastern Islands} *k}{11-20}
\begin{adjustbox}{width=\textwidth}
	\begin{threeparttable}\st{2pt}
	\begin{tabular}{llllllllll}\lsptoprule
*gloss	&	wood	&	louse	&	\tsc{1pl.excl}	&	\tsc{1pl.incl}	&	\tsc{2pl}	&	fear	&	ascend	&	root	&	fish\\
PMP	&	*\tbr{k}ahiw	&	*\tbr{k}utu	&	*\tbr{k}ami	&	*\tbr{k}ita	&	*\tbr{k}amuyu	&	*ma-ta\tbr{k}ut	&	*sa\tbr{k}ay	&	*wa\tbr{k}aR	&	*hi\tbr{k}an\\	\midrule
Bati	&	‹\tbr{k}aya›	&	‹\tbr{k}utim›	&		&		&		&		&		&	‹a\tbr{k}ar›\su{†}	&	‹i\tbr{k}an›\\
Geser	&	{\hr}\tbr{k}ai	&	{\hr}\tbr{k}utu	&	{\hr}\tbr{k}am	&	{\hr}\tbr{k}ita	&	{\hr}\tbr{k}umu	&	{\hr}mata\tbr{k}ut	&	{\hr}sa\tbr{k}a	&	{\hr}a\tbr{k}ar	&	{\hr}i\tbr{k}an\\
Gorom	&	{\hr}ai	&	{\hr}utu	&	{\hr}ami	&	{\hr}ita	&	{\hr}umu	&	{\hr}mata\tbr{ʔ}ut	&	{\hr}sa\tbr{ʔ}a	&	{\hr}a\tbr{ʔ}ar	&	{\hr}i\tbr{ʔ}an\\
Watub.	&	{\hr}\tbr{k}ai	&	{\hr}\tbr{k}utu	&	{\hr}\tbr{k}ami	&	{\hr}\tbr{k}ita	&	{\hr}\tbr{k}emiu	&	{\hr}mata\tbr{k}ut	&	{\hr}ha\tbr{k}a	&	{\hr}a\tbr{k}al	&	{\hr}i\tbr{k}an\\
Kesui	&	‹aʹi›	&	‹utun›	&		&		&		&		&		&	‹ai á\tbr{h}a›	&	‹iʹan›\\
Kowiai	&	{\hr}ai	&	{\hr}ut	&	{\hr}am	&	{\hr}ita	&	{\hr}ou	&	{\hr}-matatuʔ	&	{\hr}-sa\tbr{ʔ}a	&	\cg{(rarin)}	&	\cg{(don)}\\
		\lspbottomrule
	\end{tabular}
		\begin{tablenotes}\tnsize
			\item[†] {\citet[127]{LoskiLoski1989} give Bati \it{aʔal} ‘root’.}
		\end{tablenotes}
	\end{threeparttable}
\end{adjustbox}
\end{table}

Watubela even preserves *q/*k > \it{k} word finally.
Examples in the available data are given in \qf{11-01}.
We have so far identified only one example in which *q is not
preserved as \it{k} in Watubela: *qapəju > \it{feli-n} ‘gall’.

\noindent{\wg\st{6pt}\tblr{>{\hspace{-\tabcolsep}}l<{\hspace{-\tabcolsep}}
>{\hspace{-\tabcolsep}\enspace}
lll}
\lsptoprule
%\tbx{11-01}	&	PMP	&	Watubela	&	local gl.	\\	\midrule
%\endfirsthead
\tbx{11-01}	&	PMP	&	Watubela	&	local gl.	\\	\midrule
%\endhead
	&	*bunu\tbr{q}	&	bunu\tbr{k}	&	kill	\\
	&	*pana\tbr{q}	&	fana\tbr{k}	&	shoot (bow)	\\
	&	*daRa\tbr{q}	&	lala\tbr{k}	&	blood	\\
	&	*Ruma\tbr{q}	&	luma\tbr{k}	&	house	\\
	&	*lawa\tbr{q}	&	laa\tbr{k}	&	spider	\\
	&	*mama\tbr{q}	&	mama\tbr{k}	&	chew	\\
	&	*ma-puti\tbr{q}	&	mfuti\tbr{k}	&	white	\\
	&	*ana\tbr{k}	&	an{\tl}ana\tbr{k}	&	child	\\
	&	*manu\tbr{k}	&	manu\tbr{k}	&	bird	\\
\lspbottomrule
\end{tabular}\mbox{}\\}
